% ============================================================================
%  CAPÍTULO 5 — CONCLUSIONES
% ============================================================================

\chapter{Conclusiones}
\label{cap:conclusiones}

Esta investigación caracterizó funcionalmente a los hospitales públicos chilenos
de alta y mediana complejidad a partir de registros administrativos GRD del
período 2019--2024. El pipeline reproducible construido transforma los egresos
registrados con fines administrativos en una representación institucional basada
en producción, complejidad, diversidad, casuística, modalidad de ingreso y alta,
composición demográfica y actividad quirúrgico-obstétrica.

La respuesta a la pregunta de investigación es afirmativa: los datos GRD permiten
identificar perfiles funcionales distinguibles que no coinciden de forma estrecha
con la clasificación administrativa vigente. La solución Ward con $K=4$, aplicada
sobre una matriz de $65 \times 35$ variables, distingue tres hospitales
pediátricos, un núcleo generalista de 54 establecimientos, un grupo de cinco
hospitales con carga clínica y estancia prolongada, y tres institutos
monográficos. La elección de cuatro grupos se sostiene como un compromiso
interpretativo: la solución $K=2$ maximiza las métricas internas, pero reduce la
red a un bloque de 62 hospitales y tres casos singulares, lo que no permite el
perfilamiento funcional propuesto.

La matriz final integra 27 variables clínico-operativas directas y ocho
componentes principales de casuística. Estas componentes condensan
aproximadamente el $72{,}3\,\%$ de la varianza de las proporciones construidas a
partir de capítulos CIE-10, secciones CIE-9-MC y GRD frecuentes. La cobertura
ampliada de los campos diagnósticos y de procedimientos permite incorporar una
parte sustantiva de la información clínica registrada. Sin embargo, indicadores
como las comorbilidades promedio deben interpretarse como carga clínica
registrada, pues pueden reflejar simultáneamente diferencias de casuística y
diferencias en la exhaustividad de la codificación entre establecimientos.

La comparación con MINSAL confirma que las dos clasificaciones describen
dimensiones diferentes de la red. La concordancia entre Ward $K=4$ y las
categorías de Alta y Mediana Complejidad fue baja
(ARI $=0{,}112$; NMI $=0{,}107$). En particular, el núcleo generalista contiene
47 hospitales de Alta y 7 de Mediana Complejidad, mientras que el grupo de carga
clínica reúne 2 y 3 establecimientos, respectivamente. Esta divergencia no
implica que una clasificación sea intrínsecamente correcta y la otra incorrecta;
indica que la clasificación administrativa y el patrón de producción clínica no
son intercambiables.

La hipótesis de mayor homogeneidad funcional se evaluó directamente sobre la
misma matriz preprocesada. Ward $K=4$ obtuvo una dispersión intragrupo
normalizada WSS/TSS $=0{,}475$ y Silhouette $=0{,}401$, frente a WSS/TSS
$=0{,}944$ y Silhouette $=0{,}070$ para MINSAL. Al controlar el número de
grupos, Ward $K=2$ mantuvo menor dispersión que MINSAL
(WSS/TSS $=0{,}707$ frente a $0{,}944$). Las pruebas de permutación con 2.000
reasignaciones preservando los tamaños de grupo confirman que ambas particiones
son más homogéneas que asignaciones aleatorias comparables, aunque la distancia
respecto de la distribución nula es considerablemente mayor para las soluciones
Ward. Por tanto, la evidencia respalda que la segmentación funcional produce
grupos más homogéneos en el espacio de variables utilizado, sin que ello
constituya una validación externa de su utilidad operativa.

Los análisis de estabilidad y sensibilidad matizan el alcance de esta conclusión.
En 100 submuestras sin reposición del $80\,\%$ de los hospitales, la mediana del
ARI respecto de la solución completa fue $0{,}958$ y la del NMI fue $0{,}917$.
La transformación composicional CLR conservó exactamente la partición principal
(ARI $=1{,}000$), con un Silhouette de $0{,}438$. En cambio, la exclusión de los
años 2020--2021 produjo una concordancia menor
(ARI $=0{,}809$; NMI $=0{,}750$). En consecuencia, la separación global de los
perfiles se sostiene bajo estas perturbaciones, pero las fronteras entre grupos
pequeños y la interpretación de componentes finas de casuística son sensibles al
período y a la representación utilizada.

El análisis descriptivo posterior al agrupamiento, realizado variable por variable,
identificó diferencias significativas, tras la corrección de
Benjamini--Hochberg, en 26 de las 28 variables directas entre los cuatro grupos
de Nivel 1. Este resultado permite caracterizar los perfiles obtenidos, pero no
constituye una prueba independiente de la partición, porque las mismas variables
participaron en su construcción. La prueba de permutación del flujo completo
acota esta circularidad: ninguna de las 500 matrices sin estructura real alcanzó
las 26 variables significativas observadas. Aun así, esta evidencia no reemplaza
una validación externa con nuevos años, hospitales o variables no utilizadas para
descubrir la partición.

La interpretación de los grupos debe mantener una distinción fundamental: los
perfiles identificados describen actividad clínica registrada, pero no miden
calidad, desempeño ni eficiencia. El estudio no incorporó costos, camas,
equipamiento, dotación de personal, tiempos de espera ni desenlaces posteriores
al alta. Por ello, los resultados pueden apoyar una caracterización
complementaria de la red, pero no justifican por sí solos decisiones de
financiamiento, asignación de recursos o evaluación de desempeño institucional.

Las principales limitaciones se relacionan con la naturaleza administrativa de
los datos, la heterogeneidad en las prácticas de codificación, la inclusión de
los años afectados por la pandemia de COVID-19, la geometría composicional de
los vectores de casuística y la ausencia de validación externa independiente.
La prueba CLR indica que la clausura composicional no alteró la identidad de los
cuatro grupos en este análisis, pero no excluye que otras transformaciones, como
ILR, aporten matices adicionales. Del mismo modo, la estabilidad temporal permite
sostener la existencia de perfiles generales, pero no interpretar de manera
rígida las asignaciones que se encuentran próximas a los límites entre grupos.

Como trabajo futuro, resulta prioritario validar longitudinalmente el pipeline
con nuevos años de egresos y evaluar su generalización al incorporar hospitales
recientemente integrados al sistema GRD. También sería relevante combinar el
perfil funcional con demanda territorial, recursos humanos, camas, costos y
resultados clínicos. Esa integración permitiría pasar desde una tipología
descriptiva de la producción hacia análisis más completos de organización,
capacidad y eficiencia de la red asistencial.
